\chapter{Mixed elements: p1 and p2 in one mesh, and the minimum rule}
\label{ch:A4_mixed_elements}
\noindent\texttt{tutorials/A\_foundations/A4\_mixed\_elements.py}\\[0.75em]
\begin{outcome}
You can build a mesh where different elements carry
different polynomial orders, you know what the MINIMUM RULE does at a
p2/p1 interface (the p2 trace is constrained down to p1 — a hanging-node-
style constraint), and you can predict the consequence: global accuracy is
set by where you SPEND the p2, not how much of it you buy.
\end{outcome}

Local p-refinement is the cheapest accuracy money in FEM — IF
the high-order region covers where the solution is rich. We solve the A1
problem on a uniform mesh where only the left half (x < 0.5) is p2. The
manufactured solution comes in two flavors:
  case L: u* rich on the LEFT  (sin(2 pi x) decaying in x) — p2 well spent;
  case R: the mirror image, rich on the RIGHT — p2 wasted, p1 limits you.
The interface constraints are exactly the machinery that the Neumann band
finding (m1a findings 4b) is about: a p2 element next to a p1 element does
NOT get to keep its quadratic face modes.

\begin{expected}
\begin{verbatim}
rich-L: 4.31e-3 / 1.30e-3 / 3.56e-4   orders 1.73 / 1.87
    rich-R: 6.52e-3 / 1.64e-3 / 4.11e-4   orders 1.99 / 2.00
  (i) Both cases converge at order ~2: the p1 half sets the RATE — a chain
      is as slow as its weakest link. p-refinement buys ORDER only if it
      covers the whole domain.
 (ii) Spending p2 on the rich half helps (rich-L beats rich-R at every
      level) — but only by a CONSTANT (1.5x at level 4), and the margin
      SHRINKS under refinement. Why? Decompose the error by half (we did:
      at strong oscillation the rich half's error is ~the same under p1
      and p2!) — the minimum-rule interface constraints hold the p2 side
      to p1 accuracy in a band along the interface, and the interface here
      sits right where case L's field is still rich. The p2 payoff needs
      the p2/p1 interface pushed away from the rich region — which is
      EXACTLY the m1a Neumann-band finding (4b) wearing steady-state
      clothes.
\end{verbatim}
\end{expected}

\section*{The script}
\begin{lstlisting}
import numpy as np
from scipy.sparse.linalg import splu
import warp as wp

from diffsim.octree.build import build_uniform
from diffsim.mesh.nodes import build_mesh
from diffsim.mesh.constraints import build_constraints
from diffsim.mesh.basis import basis_tables
from diffsim.assembly.operators import DeviceMesh, assemble_csr
from diffsim.physics.poisson import (gauss_points, make_load_kernel,
                                     l2_error_masked)

DEVICE = "cuda:0"


def fields(rich_side):
    """u* rich (high curvature) on one side, gentle on the other."""
    s = +1.0 if rich_side == "L" else -1.0

    def u(x):
        # arg 2 pi (1-t)^2 has its steep gradients where t -> 0: choose t
        # so that happens on the requested side (t = x makes x=0 rich)
        t = x[:, 0] if rich_side == "L" else 1.0 - x[:, 0]
        return np.sin(1.5 * np.pi * (1 - t) ** 2) * np.sin(np.pi * x[:, 1])

    # manufactured f by high-accuracy central differences (a legitimate MMS
    # shortcut when hand-derivatives get tedious; error ~1e-10 << h^3)
    def f(x, eps=1e-5):
        lap = np.zeros(len(x))
        for c in (0, 1):
            xp = x.copy(); xp[:, c] += eps
            xm = x.copy(); xm[:, c] -= eps
            lap += (u(xp) - 2 * u(x) + u(xm)) / eps ** 2
        return -lap
    return u, f


def solve(level, rich_side):
    tree = build_uniform(level, dim=2)
    anchors = tree.anchors() / (tree.anchors().max() + (
        tree.anchors().max() == 0))
    # p2 on the LEFT half, p1 on the right (uniform level => one-knob safe)
    G = 1 << level
    xs = tree.anchors()[:, 0] / float(G)
    p_elem = np.where(xs < 0.5, 2, 1).astype(np.int8)
    mesh = build_mesh(tree, p=p_elem)
    cons = build_constraints(mesh)
    tables = {pv: basis_tables(pv, dim=2) for pv in mesh.bins}
    dm = DeviceMesh.from_mesh(mesh, cons, tables, DEVICE)
    u_star, f_star = fields(rich_side)
    A = assemble_csr(dm)
    F_full = wp.zeros(dm.n_nodes, dtype=wp.float64, device=DEVICE)
    xq = gauss_points(mesh, dm.tables_by_p)
    for pv, b in dm.bins.items():
        fq = wp.array(f_star(xq[pv]), dtype=wp.float64, device=DEVICE)
        lk = make_load_kernel(b["nbf"], b["nqp"], dm.dim)
        wp.launch(lk, dim=len(b["eids"]),
                  inputs=[b["conn"], b["h"], b["N"], b["w"], fq, F_full],
                  device=DEVICE)
    b_vec = np.asarray(cons.T.T @ F_full.numpy())
    coords = mesh.node_coords[cons.free_nodes]
    bdry = np.where(mesh.boundary_nodes[cons.free_nodes])[0]
    A = A.tolil()
    for i in bdry:
        A.rows[i] = [int(i)]
        A.data[i] = [1.0]
        b_vec[i] = u_star(coords[i:i + 1])[0]
    u_free = splu(A.tocsr().tocsc()).solve(b_vec)
    u_all = np.asarray(cons.T @ u_free)
    err = l2_error_masked(dm, u_all, u_star, lambda x: np.ones(len(x), bool))
    return err, dm.constraints.T.shape[1]


if __name__ == "__main__":
    for side in ("L", "R"):
        errs, dofs = zip(*[solve(lv, side) for lv in (4, 5, 6)])
        orders = [np.log2(errs[i] / errs[i + 1]) for i in range(2)]
        print(f"rich-{side} (p2 on left): errors "
              + "  ".join(f"{e:.3e}" for e in errs)
              + "   orders " + "  ".join(f"{o:.2f}" for o in orders)
              + f"   (dofs {dofs[-1]})")
\end{lstlisting}

\begin{explore}
\begin{verbatim}
(a) Print len(cons.rows)-style constraint counts (inspect
      build_constraints output) and find the interface constraints. How
      many are there per interface element at p2/p1? Draw the constrained
      edge mode.
  (b) Move the p2 region to a strip |x-0.5| < 0.25 straddling the
      interface of richness. Does half the p2 budget buy most of the
      benefit?
  (c) Connect to the Neumann band finding (m1a findings 4b): there, the
      constrained trace modes sat next to elements whose HESSIANS entered
      the boundary condition. Re-read the finding and restate it in this
      chapter's vocabulary.
  (d) PERFORMANCE CORNER: p2 elements cost ~ (3/2)^d more DOFs and their
      element kernels ~ (9/4)^d more work per element in 2-D. Measure
      assembly time p1-uniform vs p2-uniform vs this 50/50 mesh at level 6.
      Is mixed-p assembly the mean of the two, and why would per-BIN kernel
      launches make it so?
\end{verbatim}
\end{explore}
